\subsection{Data construction detail}
\label{app:data}

Section~\ref{sec:data} describes the two-stage synthetic-population pipeline,
parameterised by the registration threshold, and its calibration to HMRC and ONS
aggregates. This appendix records the implementation detail deferred from that
account.

Each of the approximately $5.80$ million firm rows carries a turnover, an input
expenditure, a net VAT liability ($v_i = 0.20\,(y_i - x_i)$), a calibration
weight, a frame indicator, an unregistered-stratum indicator, a VAT-scope
flag, and a registration flag. The
weights are chosen so the weighted frame rows reproduce the ONS enterprise
total and employment bands while the weighted registered rows reproduce the
HMRC trader counts and liability totals (Table~\ref{tab:calibration}). Because
the population is calibrated \emph{to} the HMRC aggregates, agreement between
my weighted outputs and those aggregates is a check of \emph{internal
consistency}, not external validation; the turnover-band rows in particular
hold by construction through the registration propensity. The headline
calibration accuracies (98.3\% and 98.2\% for the two
vintages) and the anchor comparison of Section~\ref{sec:static} are therefore
consistency checks on the pipeline, not out-of-sample tests of firm behaviour.

\begin{table}[htbp]
\centering
\caption{Calibration accuracy by dimension, $\max\{0,1-\lvert\widehat{T}-T\rvert/\lvert T\rvert\}$
averaged over the dimension's targets. HMRC dimensions are scored on
VAT-registered rows, ONS dimensions on frame rows. Source:
\texttt{results/calibration\_accuracy.txt}.}
\label{tab:calibration}
\begin{tabular}{lrr}
\toprule
Dimension (universe) & 2023--24 & 2024--25 \\
\midrule
HMRC turnover bands (registered) & 100.0\% & 100.0\% \\
ONS enterprise total (frame) & 99.8\% & 99.9\% \\
ONS employment bands (frame) & 99.9\% & 99.9\% \\
HMRC sector counts (registered) & 93.3\% & 92.6\% \\
HMRC net liability by band (registered) & 98.4\% & 98.9\% \\
\midrule
\textbf{Overall (five calibrated dimensions)} & 98.3\% & 98.2\% \\
\midrule
\emph{Diagnostic:} net liability by sector (registered) & 43.1\% & 44.0\% \\
\emph{Diagnostic:} net liability below threshold & 36.6\% & 0.0\% \\
\bottomrule
\end{tabular}
\end{table}

The reported calibration score is a bounded summary statistic. For each
dimension I compute
\(\max\{0,1-\lvert \widehat{T}-T\rvert/\lvert T\rvert\}\), so a dimension missed
by more than 100 percent receives zero rather than a negative score. The headline
``overall'' number is the unweighted mean of the five dimensions used as
calibration targets---HMRC turnover bands, ONS population, employment bands,
HMRC sector counts, and VAT liability by turnover band. VAT liability by sector
is excluded from that mean and reported separately as an informational
diagnostic, because the current generator does not calibrate sector-specific
input/output VAT structure.

The replication report also publishes the calibration-weight minimum,
median, upper quantiles, maximum, coefficient of variation, and Kish effective
sample size. These diagnostics matter because aggregate target agreement can
coexist with a small number of influential synthetic records. The report is
generated rather than transcribed, and the test suite verifies source-band
support, sector-conditional employment assignment, and weighted voluntary
registration.

\paragraph{Counterfactual exclusion window and polynomial degree.} The no-VAT
counterfactual density of Section~\ref{sec:bunching} is fitted by polynomial
regression on the observed density outside a manipulation window around the
threshold. The baseline excludes $[T^{*}-15\text{k}, T^{*}+15\text{k}]$ and fits
a degree-7 polynomial. Widening the window guards against contamination of the counterfactual
by the behavioural response but reduces the data available for the fit; raising
the degree improves in-sample fit but risks over-fitting the tails. The
sensitivity of the excess-mass and bunching-ratio estimates to window width and
degree, with associated uncertainty, is reported in
Appendix~\ref{app:inference}.
